\subsection{Image and Video Personalization}

\textbf{Personalized Image Generation.}
Prior work in personalized image generation has primarily focused on two technical directions: 1) identity-specific tuning and 2) tuning-free methods. Identity-specific tuning trains a text-to-image model to incorporate the identity by finetuning on a specific identity.
Textual Inversion~\citep{gal2022image} finetunes special text tokens for the new identity. DreamBooth~\citep{ruiz2023dreambooth} selects a few images from the same identity as reference as well as a special text token to represent the identity concept. LoRA techniques~\citep{hu2021lora} have been explored to tune a light-weight low-rank adapter
to accelerate the training process. HyperDreamBooth~\citep{ruiz2023hyperdreambooth} further reduces the training latency by directly predicting the initial weights of LoRA from the reference images. A major drawback of identity-specific tuning personalization methods is that the final model has parameters that are trained for and associated with a specific identity, and this process does not scale well to multiple users.

To overcome the limitations of the identity-specific tuning methods, another line of research extracts vision embeddings from the reference image and directly injects it into the diffusion process. This direction is more scalable as all users can share the same base model.
ELITE~\citep{wei2023elite} extracts vision features from the reference image and converts it to the text-embedding space through a local and a global mapping. PhotoMaker~\citep{li2023photomaker} merges the vision and text tokens and replaces the original text tokens for cross-attention.
PhotoVerse~\citep{chen2023photoverse} incorporates an image adapter and a text adapter to merge the vision and language tokens respectively. IP-Adapter-FaceID-Plus~\citep{ye2023ip} leverages face embedding and clip
vision encoder for identity preservation. InstantID~\citep{wang2024instantid} is a control-based method that adds ControlNet~\citep{zhang2023adding} to further control the pose and facial expression. MoA~\citep{ostashev2024moa} proposes a mixture-of-attention architecture to better fuse the vision reference and the text prompts. Imagine Yourself~\citep{meta24memu} proposed a full parallel model architecture,
a multi-stage finetuning strategy, and a novel synthetic paired data generation mechanism for better identity preservation, prompt alignment, and visual quality.

\textbf{Personalized Video Generation.}
While the aforementioned works have shown promising results in personalized image generation, adding personalization capability to video generation remains a challenging and unsolved problem. There are a few novel challenges in the area of personalized video generation: 1) compared to personalized image generation, personalized video generation needs to support more diverse and complex modifications on the reference image, \eg, turning the head, changing poses, and camera motion movements,
2) personalized video generation increases the expected quality threshold on expression and motion naturalness due to its temporal nature, and 3) finetuning a video model is much more costly than finetuning an image model, given the larger model and input sizes.

One direction of personalized video generation utilizes pose to control the video generation. Both GAN-based methods~\citep{chan2019everybody, yoon2021pose} and diffusion-based method~\citep{wang2024disco, wang2024disco, hu2024animate, xu2024magicanimate} have been proposed to generate videos following reference poses. These models are designed to animate the reference image towards the target motion, and are good for motions like singing and dancing. However, these models
require a pose sequence as reference, usually extracted from a real video, limiting its usage to a broader scope of scenearios. Also, these models often introduce occlusion and unnatural motion due to the non-ideal pose extraction and control.

On the other hand, preliminary works have been proposed to turn a personalized image model to a personalized video model. Magic Me~\citep{ma2024magic} uses identity-specific finetuning to inject identity into a video generation model. ID-Animator~\citep{he2024id} leverages a face adapter to extract identity information from single reference facial image for personalized video generation. DreamVideo~\citep{wei2024dreamvideo} and Still-Moving~\citep{chefer2024still} combine an identity adapter and a motion adapter for flexible video customization. CustomCrafter~\citep{wu2024customcrafter} proposes a
plug-and-play module for subject concept injection. Our work also focuses on using identity extraction and control methodology, as it is more generalizable to different users.
